% How a position reaches a learner: the same for every model, so each paper
% reads it from here, under its own section heading, and adds what is its own.
Every model of the lab reads a position the same way, so two models differ in
what they are asked and never in how a board reaches them.

\paragraph{The board.}
A position arrives as a FEN and becomes 64 tokens, one a square. The mover's
pawn, knight, bishop, rook, queen and king are tokens 1 to 6, the opponent's
are 7 to 12, and an empty square is 0. When Black is to move the board is
turned and the colours are swapped, so the side to move always plays up the
board. A network then learns one game, not a white one and a black one, and a
position it meets with White to move also teaches it the mirrored position
with Black to move.

\paragraph{Marks.}
What the task knows about particular squares travels as marks: yes-or-no
channels laid over the 64 squares, which a network reads together with the
pieces. This model has \V{task.marks.count} (Table~\ref{tab:marks}). The first
is always the square a pawn may be taken on in passing. The others name the
squares a move leaves and reaches, which is how a network is shown a move
without being told what the move does.

\paragraph{Context numbers.}
Whatever is not about a square is a number beside the board: the four castling
rights, the mover's first, and what the task knows about the game. This model
has \V{task.context.count} (Table~\ref{tab:reads}). Every learner reads them,
the trees that see no board included. A network reads them standardised by the
mean and spread of the train rows.

\paragraph{Hand features.}
Gradient-boosted trees cannot read 64 tokens, so the lab also computes what a
person would count on a board: material, mobility, checks and captures on
offer, pieces that can be taken, passed pawns and pressure on each king, and
for a line of moves what the line does. Table~\ref{tab:hand} in the appendix
lists all \V{task.hand.count}. They serve twice. They are what the reference
run reads, which makes them the bar for a network: a network that reads the
squares themselves should find at least what a few dozen counts already hold.
And a run can hand them to a network beside the board, where they join the
context numbers, which tests whether the network had found them on its own.

\begin{table}[htbp]\centering\small
\block{reads}
\caption{The context numbers: what every learner reads beside the board.}
\label{tab:reads}
\end{table}

\begin{table}[htbp]\centering\small
\block{marks}
\caption{The marks: what the task tells a network about squares. A square
is always given as the mover sees it, after the board is turned.}
\label{tab:marks}
\end{table}

\iflook
\begin{figure}[htbp]\centering
\block{encoding}
\caption{One position of the dataset twice. Left: as it was played, with an
arrow for the opponent's move that led to it (grey) and for the move the row
is about (blue). Right: as a network reads it, turned so that the mover plays
up the board, with a wash on every square a mark is set on: grey for the
opponent's move and, where the task marks a line of moves, blue for the
mover's and orange for the replies.}
\label{fig:encoding}
\end{figure}
\fi
